\section{Study design}
\label{sec:setup}

\subsection{Tasks and codebooks}
We use five tasks from three public, human-coded datasets (Table~\ref{tab:main_datasets}).
\begin{itemize}
  \item \textbf{MFRC} \citep{trager2022mfrc}: Reddit comments coded for moral foundations with the dataset's guidelines. The
  codebook is multi-label with 8 codes; 1,988 of the 2,000 sampled comments have three annotators and 12 have two.
  \item \textbf{AnnoMI} \citep{wu2022annomi,wu2023annomi}, full version: transcribed motivational-interviewing demonstrations, with mutually exclusive
  codes for client talk (change, sustain, neutral) and therapist behaviour (reflection, question, input, other). AnnoMI
  publishes labels rather than full definitions, so its codebooks combine AnnoMI's labels with MITI~4 wording
  \citep{moyers2016miti} and the AnnoMI-to-MISC mapping of \citet{sun2024bimisc}. A subset is coded by ten annotators.
  \item \textbf{TalkMoves} \citep{suresh2022talkmoves}: K-12 mathematics classroom transcripts coded with the dataset's
  manual, with mutually exclusive teacher moves (six plus none) and student moves (four plus none).
\end{itemize}
Together the tasks vary what makes coding hard: a sparse multi-label codebook against exclusive ones, self-contained
comments against utterances that depend on the previous turn, and low against moderate agreement among human coders
(Krippendorff's $\alpha$ of 0.17 on MFRC, 0.47--0.74 on AnnoMI). For dialogue tasks the coder also sees the previous turn.
Dunivin's coded passages about W.~E.~B.~Du Bois
\citep{dunivin2025scaling} serve as a replication task. We draw stratified samples of 2,000 comments (MFRC) or 1,500
segments (each dialogue task) that keep rare codes, and weight each segment by its inverse inclusion probability
\citep{horvitz1952generalization}. A hash of the cluster id (comment, transcript or lesson) assigns about 30\% of clusters to a
calibration split and the rest to a test split; because clusters differ in size, the calibration split holds 25--40\% of
segments. The reference label is the annotator majority; ties are dropped.

\begin{table}[t]
\centering
\small
\caption{Tasks, codebooks and samples. Calib and test are segments; clusters are conversations or transcripts (comments for MFRC). Dunivin is the replication task.}
\label{tab:main_datasets}
\begin{adjustbox}{max width=\linewidth}
\begin{tabular}{lllrrrrlr}
\toprule
task & unit & codebook & codes & calib & test & clusters & annot./item & human $\alpha$ \\
\midrule
MFRC & Reddit comment & multi-label & 8 & 608 & 1392 & 2000 & 3 & 0.17 \\
AnnoMI client & client utterance & exclusive & 3 & 569 & 931 & 128 & 1 (10 on a subset) & 0.47 \\
AnnoMI therapist & therapist utterance & exclusive & 4 & 602 & 898 & 132 & 1 (10 on a subset) & 0.74 \\
TalkMoves teacher & teacher sentence & exclusive & 7 & 370 & 1130 & 444 & 1 & -- \\
TalkMoves student & student sentence & exclusive & 5 & 383 & 1117 & 436 & 1 & -- \\
Dunivin & historical passage & multi-label & 9 & 34 & 86 & 120 & 1 & -- \\
\bottomrule
\end{tabular}
\end{adjustbox}
\par\smallskip\footnotesize\raggedright Multi-label: one yes/no decision per code; exclusive: one choice among the codes (counts include any 'none' or 'other' option). Dialogue units include the previous turn as context. $\alpha$: Krippendorff's alpha on multiply-coded items (macro over codes for multi-label); -- where items are single-coded. RQ2 uses MFRC and AnnoMI's ten-annotator subset.
\end{table}


\subsection{Coders and confidence}
Table~\ref{tab:main_coders} lists the models. \textbf{Decision models} receive identical requests: the segment and its
context, and either one multiple-choice question (exclusive codebooks) or one yes/no question per code (multi-label), whose
criteria are the code's definition, inclusion and exclusion rules and examples. OpenAI Decisions receives the same content
mapped to its own choice and predicate types; Clef, pplx-decider and Kev accepted Jev's requests unchanged. We self-host Kev
on Modal GPUs from the authors' release. \textbf{LLMs} receive the same codebook text in one prompt and return JSON with a
one-sentence rationale, an answer and a probability from 0 to 100 (per code for multi-label tasks). Appendix~\ref{app:prompts}
shows both requests.

\paragraph{From outputs to confidence.} For decision models, an exclusive decision is the most probable option and its
confidence is that probability; a multi-label decision says a code applies if its probability is at least 0.5, with
confidence $\max(p, 1-p)$. LLM signals are derived as follows.
\begin{itemize}
  \item \emph{Verbalized, exclusive:} the chosen option gets the stated probability (at least $1/K$), and the other options
  share the remainder equally. RQ2's model entropy for these signals therefore comes from a reconstructed, not an elicited,
  distribution.
  \item \emph{Verbalized, multi-label:} we read the stated probability as confidence in the stated yes/no answer. Read as
  ``probability that the code applies'', it contradicted the answer for 85\% of Luna's code decisions (0.2--4.6\% for the
  other LLMs). Confidence is therefore $c=\max(q, 1-q)$, and $p(\text{applies})$ is $c$ for ``yes'' and $1-c$ for ``no''.
  \item \emph{Token probability:} the probability of the answer token, read after the generated rationale: the option letter
  for exclusive codebooks and the yes/no token for each multi-label code (Luna, gpt-oss and both Qwen models).
  \item \emph{3-LLM agreement:} the vote share of the Luna, Haiku and gpt-oss verbal runs; \emph{self-consistency:} the vote
  share of five Luna samples at temperature 1 on 350 segments per task.
\end{itemize}
We run reasoning models with minimal reasoning so that answer-token probabilities exist. Sonnet and Qwen3.8-2.4T are the
frontier references; Qwen3.8-2.4T runs on a 270-segment-per-task subset. Every model was called on 10--12 October 2026 (UTC),
several of them in public beta (Table~\ref{tab:tab18_model_snapshot} lists the versions).

\begin{table}[t]
\centering
\small
\caption{Models and confidence signals, with realized API cost per 1,000 segments (October 2026 prices).}
\label{tab:main_coders}
\begin{adjustbox}{max width=\linewidth}
\begin{tabular}{llllr}
\toprule
model & kind & \makecell[bl]{confidence \\{} signals} & sample & \makecell[br]{USD / 1k \\{} segments} \\
\midrule
Jev & decision model & probability & full & 0.047 \\
OpenAI Decisions & decision model & probability & full & 0.102 \\
Clef & decision model & probability & full & 0.213 \\
Clef-flash & decision model & probability & full & 0.080 \\
pplx-decider & decision model & probability & full & 0.019 \\
Kev-9B & open decision model & probability & full & 0.073 \\
Kev-27B & open decision model & probability & full & 0.315 \\
Luna & LLM & verbalized; token prob.; self-consistency (350/task, \$0.28) & full & 0.107 \\
Haiku & LLM & verbalized & full & 1.906 \\
gpt-oss & LLM & verbalized; token prob. & full & 0.232 \\
Qwen3.5-9B & LLM & verbalized; token prob. & full & 0.176 \\
Sonnet & LLM (frontier) & verbalized & full & 1.964 \\
Qwen3.8-2.4T & LLM (frontier) & verbalized; token prob. & 270/task & 2.481 \\
3-LLM agreement & ensemble & vote share & full & 2.245 \\
\bottomrule
\end{tabular}
\end{adjustbox}
\par\smallskip\footnotesize\raggedright Kev: estimated serving cost from observed throughput on Modal (Kev-9B 30 calls/s on 2 H100s at \$3.95/h; Kev-27B 11 calls/s on 2 B200s at \$6.25/h), excluding cold starts and idle time. Model identifiers, settings, versions returned and call dates are in Table~\ref{tab:tab18_model_snapshot} (appendix).
\end{table}


\subsection{Decisions and metrics}
A \emph{decision} is one label to accept or reject: one per segment for exclusive codebooks, one per segment and code for
multi-label ones. \textbf{Error AUROC} (RQ1, primary) is the probability that a correct decision has higher confidence than an
incorrect one, with ties counted as one half \citep{hendrycks2017baseline}: 0.5 is chance, and 1 means every error has lower
confidence than every correct decision. \textbf{Coverage@5\%} (RQ1, primary) is the share of
test decisions accepted by the threshold rule below, reported with the error among them. Secondary metrics are the area under
the risk--coverage curve (AURC), expected calibration error (ECE, 10 bins), Cohen's $\kappa$ \citep{cohen1960coefficient} and
cost. RQ1 and EQ1--EQ5 use design weights, as do RQ3's changes; RQ2's correlations and leave-one-annotator-out $\kappa$ are
unweighted. Pooled estimates are equal-weight means over the five tasks; this describes these five tasks, not a population
of codebooks. Confidence intervals come from a cluster bootstrap that resamples calibration and test clusters separately and
refits thresholds in every replicate (1,000 replicates), unless stated otherwise.

\subsection{Threshold rules}
\label{sec:rules}
\begin{figure}[tb]
  \centering
  \begin{tikzpicture}[node distance=4mm and 7mm, every node/.style={font=\footnotesize, align=center},
      box/.style={draw, rounded corners=2pt, minimum height=8mm, inner sep=3pt}, ->, >=Stealth]
    \node[box] (corpus) {corpus of\\sampling units};
    \node[box, right=of corpus, yshift=6mm] (cal) {random sample:\\human labels\\+ model scores};
    \node[box, right=of corpus, yshift=-6mm] (rest) {remainder,\\uncoded};
    \node[box, right=of cal] (fit) {threshold $\hat t$\\(naive or LTT)};
    \node[box, right=of rest] (model) {model: label +\\confidence $c$};
    \node[box, right=of model, xshift=3mm] (gate) {$c \ge \hat t$?};
    \node[box, right=of gate, yshift=6mm] (acc) {accept label};
    \node[box, right=of gate, yshift=-6mm] (hum) {human codes it};
    \draw (corpus) -- (cal); \draw (corpus) -- (rest); \draw (rest) -- (model); \draw (model) -- (gate);
    \draw (cal) -- (fit); \draw (fit) -| (gate); \draw (gate) -- node[above, sloped] {yes} (acc); \draw (gate) -- node[below, sloped] {no} (hum);
  \end{tikzpicture}
  \caption{The triage workflow we evaluate. A random sample of the corpus's units (here comments, transcripts or lessons) is
  coded by hand and by the model; comparing the two fixes a threshold. The model's confidence then decides which labels in the
  rest of the corpus are accepted and which go to a person. A guaranteed rule bounds the error rate among the accepted labels
  of the remainder.}
  \label{fig:workflow}
\end{figure}
A researcher who has hand-coded a sample needs a rule that turns it into a threshold. Each rule below picks a confidence
threshold $t$ from the hand-coded calibration clusters and accepts the other decisions with confidence $\ge t$; they range
from common practice to a rule with a guarantee. We distinguish two targets. The \emph{corpus target} is the design-weighted error among the accepted
decisions of the uncoded remainder of the corpus, the quantity a researcher coding their own data cares about. The
\emph{population target} is the expected error among accepted decisions for new clusters from the same source. Whether one
particular test set has error above 5\% is exactly the corpus target when that test set is the uncoded remainder of a
random split, and only a diagnostic for the population target.
\begin{itemize}
  \item \textbf{Naive} (the pre-specified primary rule): the lowest threshold whose calibration error is at most 5\%. No guarantee.
  \item \textbf{LTT on decisions:} Learn-then-Test \citep{angelopoulos2025ltt} with binomial tests of ``risk $>$5\%'' on
  decision counts, fixed-sequence from the most confident tie block. With design weights the counts are Kish-effective and
  fractional, so the test is an approximation even for independent decisions. It treats decisions as independent; we show it
  only for comparison.
  \item \textbf{Corpus LTT} (our guaranteed rule). Intuitively, the test gathers evidence, cluster by cluster, that the
  accepted decisions are wrong less than 5\% of the time, and certifies a threshold only when that evidence is strong enough
  to extend to the rest of the corpus. Formally, for each threshold on a grid of 16 values from 0.5 to 1.0, cluster $i$
  contributes $Z_i(t)=\sum_j w_{ij}\,\mathbf{1}\{c_{ij}\ge t\}\,(e_{ij}-0.05)$, so a set of clusters has error at most 5\% among
  its accepted decisions exactly when its $Z$-sum is at most 0. Rescaled by $B$, the largest total weight of one cluster in the
  corpus (known before anything is coded), $Y_i=(Z_i+0.05B)/B$ lies in $[0,1]$. A betting test for sampling without
  replacement \citep{waudbysmith2024betting} visits the hand-coded clusters in random order and asks whether the corpus mean
  of $Y$ is low enough for the remainder's $Z$-sum to be at most 0 (Appendix~\ref{app:ltt}). So that the result does not
  depend on one arbitrary order, we average the test's final wealth over 100 random orders. Each threshold is tested at
  $0.05/16$, and the lowest certified one is used.
\end{itemize}
\noindent\textbf{Proposition} (corpus LTT; Appendix~\ref{app:ltt}). Suppose the hand-coded clusters are a uniformly random
subset of the corpus's clusters, the orders in which the test visits them are uniformly random and independent of the labels,
and the confidence scores, the threshold grid and $B$ are fixed before the hand-coded labels are used. Then, with probability
at least 0.95 over the random choice of clusters and orders, every certified threshold has error at most 5\% among the
accepted decisions of the uncoded clusters (counting a remainder with no accepted decision as within the target). No
assumption is made about how errors depend on each other within or across clusters.

The construction adapts existing testing machinery to the coding workflow; the statistical tools are not new. Nothing in
it requires the unit to be a cluster: any unit can be sampled, such as a comment, a whole transcript, a fixed-size block of
turns or a single utterance coded with its context, provided the hand-coded units are a random sample of the corpus's units
and $B$ is the largest unit's weight. What the guarantee must respect is how the hand-coded sample was drawn.

\paragraph{How we apply it.} In our evaluation, each task's sample is the corpus, its calibration split the hand-coded
clusters (assigned by a hash of the cluster id, which we treat as random) and its test split the remainder. The visiting
orders come from a pseudo-random generator with a seed recorded before this analysis was run. The guarantee holds for each
configuration on its own; picking the best of several certified configurations would need a further correction. Learned
scores, such as per-code maps or the prevalence baseline, are fitted on half of the calibration clusters, frozen, certified
on the other half and applied unchanged to the remainder. Multi-stage cascades are certified jointly over all combinations
of per-stage thresholds.

\paragraph{Two reference rules.} For the population target we also report the same test with a constant null mean
(\emph{LTT on i.i.d.\ clusters}). It needs clusters to be independent draws, which we assume rather than know, and $B$ to
bound every future cluster's weight, which is true for MFRC (eight codes per comment at one weight) but not for transcripts
of unbounded length. The \textbf{oracle} applies the naive rule to the test split itself: it is not deployable, and serves
as an upper bound for comparison.

\subsection{Analysis plan and reference labels}
\label{sec:plan}
\paragraph{Analysis plan.} The plan fixed the tasks, splits, decision units, metrics, primary rule and RQ1 contrasts before
any test-split outcome was computed. Four amendments, each dated before the data it governs, added OpenAI Decisions and a
Qwen comparator (v1.1), five more decision models (v1.2), a memorization probe (v1.3) and a prompt-sensitivity check (v1.4).
The plan and amendments were kept in the project repository, not registered with an external service, so their timing
cannot be verified independently; we call the plan prospectively specified. A claim needs the pooled 95\% CI to exclude
zero; for the v1.2 family the plan also requires Holm's correction \citep{holm1979simple}, which we apply to the earlier
families as a sensitivity analysis. Jev anchors the confirmatory contrasts because the original plan was written for it;
the other decision models were added by amendment. Pooled intervals resample each corpus's clusters jointly, so AnnoMI's two
tasks share transcript draws and TalkMoves' two share lesson draws. Sonnet was specified on a nested subset; that result is the
confirmatory one, and its later full-sample run is a labelled extension. Everything changed after results were seen is
labelled exploratory and logged (Appendix~\ref{app:deviations}); this includes the corpus guarantee, which replaced
decision-level LTT and was revised after external reviews.

\paragraph{Adjudicated reference.} Human majority labels can contain errors. We sampled 200 items for blinded adjudication from
disagreements between the human label and six label sources (Jev in both question formats, Luna, Haiku, gpt-oss and Sonnet's
subset run), oversampling high-confidence disagreements, plus hidden controls. A two-family LLM panel (Claude Opus 5.5
and GPT-6 Sol) judged each item. The author adjudicated 36 of the 52 items the panel disputed, plus a 12-item audit of panel-agreed items (agreement 10 of 12); a planned
second sitting to measure the author's own consistency was not done. Of the 200 items, 184 were resolved; unresolved items
keep the human label. We use this reference as a sensitivity analysis, not as a corrected gold standard.
